% Part: proof-theory
% Chapter: proof-search
% Section: search-algorithm

\documentclass[../../../include/open-logic-section]{subfiles}

\begin{document}

\olfileid{pt}{ps}{sal}

\olsection{خوارزمية البحث في~\Log{G3c}}

نصف خوارزمية بحث للنسق~\Log{G3c}. وليست هذه الخوارزمية الوحيدة الممكنة،
ولا حتى أكفأها، لكنها صحيحة: فإذا كان للمتتالية~$\Gamma\Sequent\Delta$
!!a{proof}، توقفت في النهاية ووجدت~!!a{proof}. وهي مضمونة أيضًا: فإذا
توقفت من دون أن تجد~!!a{proof}، أو لم تتوقف قط، كانت
$\Gamma\Sequent\Delta$ غير صالحة أصلًا.

من الخصائص الأساسية اللازمة أن تضمن الخوارزمية ``اختزال'' كل~!!{formula}
في~$\Gamma\Sequent\Delta$ أو في أي متتالية تنتج أثناء البحث؛ أي معاملتها
بوصفها الصيغة الرئيسة لقاعدة استدلال مطبقة إلى الخلف، وأن يتكرر ذلك كلما
لزم. ونسمي هذه الخاصية \emph{الإنصاف}.

تعمل الخوارزمية على مراحل. وناتج كل مرحلة شجرة متتاليات، ويُنتج ناتج المرحلة
التالية باختزال~!!a{formula} في كل متتالية عليا ليست مسلّمة بالفعل.

في المرحلة~$0$ نبدأ بـ~$\Gamma\Sequent\Delta$. ولضمان الإنصاف نضع أعدادًا
مؤشراتٍ لـ~!!{formula}s في~$\Gamma$ و$\Delta$ بحيث تتلقى~!!{formula}s
المختلفة أعدادًا مختلفة. ويمكن، مثلًا، عد~!!{formula}s من اليسار إلى
اليمين. ونكتب الأعداد مؤشرات علوية. فإذا كنا نبحث، مثلًا، عن~!!a{proof}
لـ~$!A,!B\Sequent !C$، كان ناتج المرحلة~$0$ هو
$!A^1,!B^2\Sequent !C^3$. وأكبر~$n$ يظهر مؤشرًا علويًا في متتالية مفهرسة
كهذه هو \emph{مؤشرها الأقصى} (وهو~$3$ في المثال).

إذا احتوت المتتالية أسوارًا، احتجنا أيضًا إلى معرفة الحدود التي نستعملها
عند اختزال~$\lexists[x][!A(x)]$ على اليمين أو~$\lforall[x][!A(x)]$ على
اليسار. ولن نحتاج إلا إلى حدود مبنية من~!!{constant}s~$\Obj c_0$ و
$\Obj c_1$ و$\Obj c_2$، \dots، باستعمال رموز الدوال الواقعة في~$\Gamma$
و$\Delta$. ونفترض أن مجموعة هذه الحدود كلها معطاة بتعداد ثابت، بحيث يمكننا
الحديث، مثلًا، عن ``أول~!!{constant} لا يقع في
$\Gamma\Sequent\Delta$.'' ولكل متتالية مفهرسة في الشجرة التي تولدها
الخوارزمية مجموعة مرتبطة من~!!{constant}s. ومجموعة~!!{constant}s المسندة إلى المتتالية
في المرحلة~$0$ هي مجموعة جميع~!!{constant}s الواقعة فيها إن وُجدت، وهي
$\{\Obj c_0\}$ إذا لم تحتوِ أي~!!{constant}s.

افترض أن الخوارزمية ولّدت في المرحلة~$n$ شجرة متتاليات مفهرسة، مع مجموعات
!!{constant}s المرتبطة بها. وفي المرحلة~$n+1$ نفعل ما يأتي لكل متتالية
عليا.

إذا وُجدت~!!a{formula}~$!A$ تقع على جانبي المتتالية معًا، أو كان جانبها
الأيسر يحتوي~$\lfalse$، فلا نفعل شيئًا: فلهذه المتتاليات~!!{proof}s في
\Log{G3c}، ولا يبقى شيء لنفعله لهذه المتتالية العليا بعينها.

إذا لم توجد~!!{formula} غير ذرية في المتتالية، فأجهض الخوارزمية. فلا يوجد
برهان لـ~$\Gamma\Sequent\Delta$.

وإلا فاختر~!!{formula} غير الذرية~$!A$ ذات أصغر مؤشر~$i$. وعندئذ تكون
المتتالية العليا إما~$!A^i,\Pi\Sequent\Lambda$ وإما
$\Pi\Sequent\Lambda,!A^i$. ولتكن~$k$ المؤشر الأقصى للمتتالية، ولتكن~$C$
مجموعة~!!{constant}s المسندة إليها.

``نختزل''~$!A^i$؛ أي نولّد متتاليات عليا جديدة وفق قاعدة الاستدلال التي
يحددها~!!{main operator} لـ~$!A$، ووفق وقوع~$!A^i$ على اليسار أو اليمين.

\begin{enumerate}
  \item إذا كان~$!A\ident !B\land !C$ ووقع على اليسار، فأضف متتالية عليا
  جديدة فوق~$(!B\land !C)^i,\Pi\Sequent\Lambda$:
  \[
  \Axiom$!B^{k+1}, !C^{k+2}, \Pi \fCenter \Lambda$
  \RightLabel{\LeftR{\land}}
  \UnaryInf$(!B \land !C)^i, \Pi \fCenter \Lambda$
  \DisplayProof
  \]
  ومجموعة~!!{constant}s المسندة إلى المتتالية الجديدة هي~$C$.
  \item إذا كان~$!A\ident !B\land !C$ ووقع على اليمين، فأضف متتاليتين
  عليين جديدتين فوق~$\Pi\Sequent\Lambda,(!B\land !C)^i$:
  \[
  \Axiom$\Pi \fCenter \Lambda, !B^{k+1}$
  \Axiom$\Pi \fCenter \Lambda, !C^{k+1}$
  \RightLabel{\RightR{\land}}
  \BinaryInf$\Pi \fCenter \Lambda, (!B \land !C)^i$
  \DisplayProof
  \]
  ومجموعة~!!{constant}s المسندة إلى كل متتالية جديدة هي~$C$.

  \item حالات~$!A\ident\lnot !B$ و$!A\ident !B\lor !C$ و
  $!A\ident !B\lif !C$ مماثلة، وتُترك تمارين.

  \item إذا كان~$!A\ident\lexists[x][!B(x)]$ ووقع على اليسار، فأضف
  متتالية عليا جديدة فوق~$\lexists[x][!B(x)]^i,\Pi\Sequent\Lambda$:
  \[
  \Axiom$!B(c)^{k+1}, \Pi \fCenter \Lambda$
  \RightLabel{\LeftR{\lexists}}
  \UnaryInf$\lexists[x][!B(x)]^i, \Pi \fCenter \Lambda$
  \DisplayProof
  \]
  حيث~$c$ أول~!!{constant} لا ينتمي إلى~$C$. ومجموعة~!!{constant}s
  المسندة إلى المتتالية الجديدة هي~$C\cup\{c\}$.
  \item إذا كان~$!A\ident\lexists[x][!B(x)]$ ووقع على اليمين، فأضف
  متتالية عليا جديدة فوق~$\Pi\Sequent\Lambda,\lexists[x][!B(x)]$:
  \[
  \Axiom$\Pi \fCenter \Lambda, \lexists[x][!B(x)]^{k+1}, !B(t)^{k+2}$
  \RightLabel{\RightR{\lexists}}
  \UnaryInf$\Pi \fCenter \Lambda, \lexists[x][!B(x)]^i$
  \DisplayProof
  \]
  حيث~$t$ أول حد لا يحتوي إلا~!!{constant}s في~$C$ ولم يُستعمل من قبل في
  اختزال~$\lexists[x][!B(x)]$ على اليمين تحت هذه المتتالية. وإذا استُعملت
  كل هذه الحدود، نعيد استعمال أول حد في التعداد الثابت مبني فقط من ثوابت
  $C$، بدل إدخال ثابت قد يقع خارج~$C$. ومجموعة~!!{constant}s المسندة إلى
  المتتالية الجديدة هي~$C$.
  \item حالات~$!A\ident\lforall[x][!B(x)]$ مماثلة، وتُترك تمارين.
\end{enumerate}

إذا لم نضف في المرحلة~$n+1$ متتالية جديدة إلى أي متتالية عليا، توقفت
الخوارزمية بنجاح. وفي هذه الحالة وجدنا~!!a{proof} لـ
$\Gamma\Sequent\Delta$، نحصل عليه من شجرة المرحلة~$n+1$ بحذف المؤشرات
كلها. فالمتتاليات العليا كلها إما من الصورة
$!A,\Pi\Sequent\Lambda,!A$ وإما من الصورة
$\lfalse,\Pi\Sequent\Lambda$. وجميع الاستدلالات صحيحة في~\Log{G3c}.
وهذا واضح للاستدلالات القضاياوية. ولاحظ أن مجموعة~$C$ من~!!{constant}s المسندة إلى متتالية
تحتوي، في كل خطوة، جميع~!!{constant}s الواقعة في المتتالية. ولذلك يتحقق
شرط المتغير الذاتي في الاختزال باستعمال~\LeftR{\lexists} و
\RightR{\lforall}، لأن المتغير الذاتي~$c$ كان~!!a{constant} لا ينتمي إلى
المجموعة~$C$ المسندة إلى النتيجة، فلم يقع~$c$ في النتيجة أيضًا. فنحصل على:

\begin{prop}
خوارزمية البحث عن البرهان في~\Log{G3c} صحيحة: إذا توقفت بنجاح، كان
للمتتالية الابتدائية~$\Gamma\Sequent\Delta$ !!a{proof}.
\end{prop}





\end{document}
